\section*{Part III: Hiring Educators and Instructional Specialists}

\subsection*{21. A Subject Expert Is Not Automatically an Effective Teacher}

The visible performance of teaching is a lecture; the underlying job is creating conditions in which learners build usable knowledge and skill. An instructor can be fluent, entertaining, and technically impressive while producing little independent learning. Another can be less theatrical yet diagnose misconceptions, make hard ideas comprehensible, create useful practice, and respond accurately to evidence of confusion. Hiring therefore cannot end at the demonstration lecture.

Evaluate educators on at least \textbf{subject competence, diagnostic skill, instructional explanation, practice and feedback design, learner responsiveness, and professional reliability}. Depending on the job, assessment literacy, online facilitation, learner inclusion, and curriculum development deserve separate attention. A charismatic teaching demo should never conceal an inability to explain \textit{why} an answer is wrong or to distinguish conceptual misunderstanding from a calculation error. Chapter 11 explained the mechanisms of learning; Chapter 12 designed instruction; Chapter 13 addressed outcome evidence. Hiring faculty means observing whether a candidate can enact those principles rather than merely repeat their vocabulary.

\subsection*{22. Start the Educator Assessment with Content, Not Performance Style}

A technical instructor should 1st clear a threshold of competence in the material. For a Bayesian statistics educator, that may mean explaining conjugacy, posterior predictive reasoning, model checks, uncertainty, and when a chosen model fails. For an AI engineer-instructor, it may mean demonstrating a deployed system, evaluation methodology, security judgment, and limitation awareness. A teacher who cannot reason accurately about the subject will not become trustworthy through excellent slides.

Assessments must match \textbf{what the role actually teaches}. Asking a beginner's Python instructor to prove a theorem in measure theory may be an irrelevant prestige test. Asking a senior probability instructor to explain measure-theoretic independence is entirely different. Give candidates a syllabus-level competency map and use short problems designed around central misconceptions. Score correctness, reasoning, boundary conditions, and calibration of uncertainty; do not disproportionately reward speed unless rapid response is genuinely necessary for the job.

\subsection*{23. The Demonstration Lesson Must Contain Learner Evidence}

A useful teaching audition is not ``present a topic of your choice for 20 minutes.'' That invites candidates to perform their best rehearsed explanation to an unnaturally attentive audience. Instead: disclose the target learner level; provide a real misconception and learning objective; ask for a 12-minute lesson; include an interruption or learner response; then ask the candidate to adjust. A reviewer may role-play a learner honestly, using a predetermined scenario rather than adversarial theatrics.

\textbf{Example}: A learner says, ``The model got 96\% accuracy, so it should work well in production.'' The candidate must explain what the metric does and does not establish, introduce the necessary comparison or data-split issue, ask a check-for-understanding question, and decide what the learner should practice next. Score whether the teacher notices the problem, explains without distorting, elicits reasoning, and adapts to the response. High-quality slides alone count for little.

\subsection*{24. The Teach-Diagnose-Reteach Exercise}

A stronger educator assessment has 3 phases:
\begin{enumerate}[leftmargin=*]
\item \textbf{Teach}: candidate gives a short explanation.
\item \textbf{Diagnose}: reviewer supplies an anonymized learner answer, including a non-obvious error; candidate identifies the misconception and the evidence supporting the diagnosis.
\item \textbf{Reteach}: candidate revises the explanation or devises the smallest next practice task that would test the hypothesis.
\end{enumerate}
This probes capability to improve learning rather than only to speak confidently.

For a statistics module, the learner says, ``The 95\% confidence interval means there is a 95\% chance that this particular fixed parameter lies inside this interval.'' An excellent candidate explains the frequentist coverage interpretation without pretending that the language issue is trivial; identifies the misunderstanding between a long-run procedure and a post-data probability statement; and designs a simulation or counterexample appropriate to the learner. A weak answer simply repeats ``that is incorrect.''

\subsection*{25. Evaluate Feedback as an Instructional Product}

In a rigorous educational institute, students pay in part for informed human attention. Ask instructor finalists to review \textbf{2 anonymized pieces of work}: 1 with a fundamental error, 1 technically correct but incomplete in reasoning. Give the same rubric and a realistic feedback-time budget. Score whether the candidate identifies important issues, prioritizes them, explains a next action, avoids completing the task for the learner, and knows when to escalate.

An exemplary response might say: ``Your pipeline leaks the target into a preprocessing step. Refit preprocessing inside each training fold; rerun the baseline before interpreting model improvements. Your error analysis can wait until the evaluation is valid.'' That is more valuable than 15 paragraphs of generic praise. Record feedback minutes as well as quality; an excellent review that takes 3 hours may not be economical for a high-volume course. Conversely, cheap instant feedback that misses the central defect is not a saving.

\subsection*{26. Observe Instructional Judgment Under Mixed Ability}

A cohort never consists of identical learners. Give a teaching candidate a scenario: 5 learners have completed the exercise; 4 cannot begin; 2 are applying the wrong abstraction; 1 is disengaged; 10 minutes remain. Ask how the instructor will decide what to do, communicate to the group, and protect individual learning. The answer need not follow 1 script. It should show a coherent method for diagnosing what is shared, providing a common intervention when justified, offering extension work for fast movers, and avoiding the humiliation of struggling learners.

Avoid rewarding only an elaborate verbal plan. Ask the candidate to write the next prompt, draw the 2-minute explanation, and explain what evidence would signal improvement. A candidate who says ``I would personalize everything'' without accounting for 20 learners and limited time has not solved the instructional problem.
